\section{Conclusion}

The transition of crypto mining companies into AI data centers is, in essence, a "migration of computing power." In this migration, tax policy is one of the important variables affecting project location, investment timing and business model design.

Mining companies need to consider more than which state offers the greatest incentives. They must also assess whether the project can qualify for the incentive, how long the benefit will remain available, and whether the project can continue to satisfy investment, employment, reporting and recordkeeping requirements after qualification. Filing an application too late, or making purchases before obtaining the required exemption certification, may result in the loss of tax savings that would otherwise have been available. Failure to meet performance commitments during subsequent operations may lead to repayment of tax, interest, or even loss of eligibility for the incentive.

U.S. states are currently at a subtle turning point in their approach to data centers. While some jurisdictions continue to offer substantial tax incentives, others have suspended approval of new projects or narrowed the scope of existing benefits. As policy windows contract, the "act first and apply later" approach is becoming increasingly unworkable. For crypto mining companies planning a transition, applications for tax incentives and long-term tax planning need to proceed in parallel with investment decisions.

\FloatBarrier
